\documentclass[14pt,a4paper]{extreport}
% preamble.tex — спільна преамбула для всіх лабораторних робіт. Підключається одразу після \documentclass у кожному lab*/main.tex.
%
% ЗБІРКА: LuaLaTeX (не pdfLaTeX). Причина: пакет listings під pdfLaTeX не читає UTF-8 і зупиняє збірку на кириличних коментарях у лістингах програмного коду. LuaLaTeX обробляє UTF-8 нативно.
%
% ШРИФТ: Liberation Serif — метрично сумісний з Times New Roman і містить повний набір кириличних гліфів (TeX Gyre Termes кирилиці не має, через що український текст не відображався взагалі).

\usepackage{fontspec}
\setmainfont{Liberation Serif}
\setmonofont{DejaVu Sans Mono}[Scale=0.82]

\usepackage{polyglossia}
\setmainlanguage{ukrainian}
\setotherlanguage{english}

\usepackage{microtype}

% Поля, інтервал, відступи й заголовки нижче взяті з шаблону НУЛП, за яким написана бакалаврська робота (../bachelor-thesis/thesis/main.tex). Той самий формат приймає В.М. Шиманський і для магістерських, тож усі звіти тримаємо в ньому, а не в тому, що переказували на консультації.
\usepackage[left=25mm,right=15mm,top=15mm,bottom=23mm,footskip=1.5mm]{geometry}
\usepackage{setspace}
% Саме \setstretch{1.5}, а не \onehalfspacing: останній дає близько 1.25.
\setstretch{1.5}

\usepackage{graphicx}
\usepackage{float}
\usepackage{booktabs}
\usepackage{array}
\usepackage{tabularx}
\usepackage{longtable}
\usepackage{enumitem}
% labelindent вирівнює маркер списку з абзацним відступом — так у шаблоні НУЛП.
\setlist{nosep}
\setlist[itemize]{label=\textbullet,leftmargin=*,labelindent=1.27cm}
\setlist[enumerate]{leftmargin=*,labelindent=1.27cm}
\usepackage{indentfirst}
\usepackage{amsmath}

\usepackage{hyperref}
\hypersetup{colorlinks=true,linkcolor=black,urlcolor=blue,citecolor=black}

% xurl дозволяє \url розривати адресу в будь-якому місці, а не лише після скісної риски. Без нього довгий сегмент шляху (напр. «scopus---podrobnaya- instruktsiya-dlya-uchenogo») виходить за межі сторінки, бо дефіс для hyperref не є точкою розриву. Вантажиться після hyperref — перевизначає його \UrlBreaks.
\usepackage{xurl}

% ---------------------------- Список використаної літератури ----------------------------
% biblatex + biber зі стилем IEEE: посилання «[1]» у тексті, записи — за порядком першої згадки. Обрано IEEE, а не gost-numeric (ДСТУ 8302:2015), — так вирішив автор звітів після порівняння обох варіантів.
% Кожен звіт із джерелами сам підключає свій references.bib через \addbibresource і друкує список \printbibliography[heading=none] під власним \section. Звіт без джерел від завантаженого biblatex нічого не втрачає.
% Оточення списку — стандартне для стилю: мітка від лівого поля, текст опису вирівняний після неї. Не підганяємо його під абзацний відступ нумерованих переліків тексту (labelindent у \setlist вище) — список літератури є окремим розділом зі своїм стандартним оформленням.
\usepackage[backend=biber,style=ieee]{biblatex}
% biblatex за замовчуванням додає розтяжний проміжок після «/ . : ? =» в адресах, і з вимкненими переносами рядок вирівнюється розрідженням самої адреси («https : / / www . scopus»). Рвати адресу й так дозволяє xurl, тож проміжок лише шкодить.
\setlength{\biburlbigskip}{0mu}
% Без додаткового інтервалу між описами: розділ має йти суцільним текстом, як решта звіту. Так само в шаблоні бакалаврської.
\setlength{\bibitemsep}{0pt}
\setlength{\bibparsep}{0pt}

% ---------------------------- Лістинги програмного коду ---------------------------- Використано minted (через Pygments), а не listings: listings не вміє коректно обробляти багатобайтові символи й «висмикує» кириличні коментарі з коду, зліплюючи їх в один рядок. Потребує прапорця -shell-escape.
\usepackage{xcolor}
\usepackage[newfloat]{minted}
% labelsep=period дає «Рис. 2.1. Назва» замість «Рис. 2: Назва» — формат НУЛП.
\usepackage[labelsep=period]{caption}

% baselinestretch=1 обов'язковий: без нього лістинг успадковує півторашку основного тексту, і довгий фрагмент коду перестає вміщатися у сторінку. Усередині [H]-флоата це не помилка збірки, а нескінченний цикл порожніх сторінок — lab1 так намолотив 161 тисячу сторінок.
\setminted{
  baselinestretch=1,
  fontsize=\footnotesize,
  linenos=true,
  frame=single,
  framesep=4pt,
  breaklines=true,
  breakanywhere=true,
  tabsize=4,
  numbersep=7pt,
  bgcolor=gray!8,
}
\setminted[python]{python3=true}

\SetupFloatingEnvironment{listing}{name=Лістинг,listname=Перелік лістингів}

% ---------------------------- Нумерація та оформлення заголовків ----------------------------
\usepackage{titlesec}

\setcounter{secnumdepth}{3}
\renewcommand{\thesection}{\arabic{section}}
\renewcommand{\thesubsection}{\arabic{section}.\arabic{subsection}}
\renewcommand{\thesubsubsection}{\arabic{section}.\arabic{subsection}.\arabic{subsubsection}}

% Розділ по центру, 16 пт, великими літерами; слово «РОЗДІЛ» опущене, бо звіт не дисертація. Номер лишається: \numberwithin прив'язує до нього нумерацію рисунків («Рис. 2.1»), і без нумерованих розділів вона б зламалася. \MakeUppercase четвертим аргументом отримує текст заголовка, тому в джерелі заголовки пишуться звичайним регістром.
\titleformat{\section}[block]
  {\centering\bfseries\fontsize{16}{19.2}\selectfont}
  {\thesection.}{0.3em}{\MakeUppercase}
% Підрозділи — зліва з абзацного відступу, звичайним регістром.
\titleformat{\subsection}[block]{\bfseries}{\thesubsection.}{1em}{}
\titleformat{\subsubsection}[block]{\bfseries\itshape}{\thesubsubsection.}{1em}{}

% Незіркові \titlespacing — обов'язкові. Без них titlesec застосовує зірковий варіант, який прибирає абзацний відступ у першого абзацу після заголовка й перекриває indentfirst. Значення — типові для report.
\titlespacing{\section}{0pt}{20pt}{6pt}
\titlespacing{\subsection}{1.27cm}{6pt}{0pt}
\titlespacing{\subsubsection}{1.27cm}{6pt}{0pt}

\setlength{\parindent}{1.27cm}
\setlength{\parskip}{0pt}
\renewcommand{\arraystretch}{1.3}

% Переноси слів вимкнені — так у шаблоні НУЛП. Натомість абзац розганяється міжслівними пробілами, тому tolerance і emergencystretch підняті: без них рядки виходили б за праве поле.
\hyphenpenalty=10000
\exhyphenpenalty=10000
\tolerance=9999
\setlength{\emergencystretch}{3em}

% ---------------------------- Підписи до рисунків і таблиць ----------------------------
\captionsetup[figure]{name=Рис.,justification=centering,font={bf,it}}
\captionsetup[table]{name=Таблиця,justification=raggedleft,singlelinecheck=false,labelfont=bf,textfont=normalfont}

% Рисунки, таблиці й формули нумеруються в межах розділу: «Рис. 2.1», а не «Рис. 5». Так у шаблоні НУЛП, і так само виглядали приклади на консультації.
\numberwithin{equation}{section}
\numberwithin{figure}{section}
\numberwithin{table}{section}

% ---------------------------- Колонтитули ----------------------------
\usepackage{fancyhdr}
\pagestyle{fancy}
\fancyhf{}
\fancyfoot[R]{\thepage}
\renewcommand{\headrulewidth}{0pt}
\fancypagestyle{plain}{%
  \fancyhf{}
  \fancyfoot[R]{\thepage}
  \renewcommand{\headrulewidth}{0pt}
}

% titlepage.tex — титульний аркуш, спільний для всіх дисциплін.
% Потребує визначених у subject.tex команд:
%   \SubjectName, \LecturerName, \StudentName, \StudentGroup, \AcademicYear
%
% Використання у тексті: \MakeLabTitle{1}{Тема роботи}
%
% Рід роботи. Дисципліни з практичними роботами перевизначають його після
% % preamble.tex — спільна преамбула для всіх лабораторних робіт. Підключається одразу після \documentclass у кожному lab*/main.tex.
%
% ЗБІРКА: LuaLaTeX (не pdfLaTeX). Причина: пакет listings під pdfLaTeX не читає UTF-8 і зупиняє збірку на кириличних коментарях у лістингах програмного коду. LuaLaTeX обробляє UTF-8 нативно.
%
% ШРИФТ: Liberation Serif — метрично сумісний з Times New Roman і містить повний набір кириличних гліфів (TeX Gyre Termes кирилиці не має, через що український текст не відображався взагалі).

\usepackage{fontspec}
\setmainfont{Liberation Serif}
\setmonofont{DejaVu Sans Mono}[Scale=0.82]

\usepackage{polyglossia}
\setmainlanguage{ukrainian}
\setotherlanguage{english}

\usepackage{microtype}

% Поля, інтервал, відступи й заголовки нижче взяті з шаблону НУЛП, за яким написана бакалаврська робота (../bachelor-thesis/thesis/main.tex). Той самий формат приймає В.М. Шиманський і для магістерських, тож усі звіти тримаємо в ньому, а не в тому, що переказували на консультації.
\usepackage[left=25mm,right=15mm,top=15mm,bottom=23mm,footskip=1.5mm]{geometry}
\usepackage{setspace}
% Саме \setstretch{1.5}, а не \onehalfspacing: останній дає близько 1.25.
\setstretch{1.5}

\usepackage{graphicx}
\usepackage{float}
\usepackage{booktabs}
\usepackage{array}
\usepackage{tabularx}
\usepackage{longtable}
\usepackage{enumitem}
% labelindent вирівнює маркер списку з абзацним відступом — так у шаблоні НУЛП.
\setlist{nosep}
\setlist[itemize]{label=\textbullet,leftmargin=*,labelindent=1.27cm}
\setlist[enumerate]{leftmargin=*,labelindent=1.27cm}
\usepackage{indentfirst}
\usepackage{amsmath}

\usepackage{hyperref}
\hypersetup{colorlinks=true,linkcolor=black,urlcolor=blue,citecolor=black}

% xurl дозволяє \url розривати адресу в будь-якому місці, а не лише після скісної риски. Без нього довгий сегмент шляху (напр. «scopus---podrobnaya- instruktsiya-dlya-uchenogo») виходить за межі сторінки, бо дефіс для hyperref не є точкою розриву. Вантажиться після hyperref — перевизначає його \UrlBreaks.
\usepackage{xurl}

% ---------------------------- Список використаної літератури ----------------------------
% biblatex + biber зі стилем IEEE: посилання «[1]» у тексті, записи — за порядком першої згадки. Обрано IEEE, а не gost-numeric (ДСТУ 8302:2015), — так вирішив автор звітів після порівняння обох варіантів.
% Кожен звіт із джерелами сам підключає свій references.bib через \addbibresource і друкує список \printbibliography[heading=none] під власним \section. Звіт без джерел від завантаженого biblatex нічого не втрачає.
% Оточення списку — стандартне для стилю: мітка від лівого поля, текст опису вирівняний після неї. Не підганяємо його під абзацний відступ нумерованих переліків тексту (labelindent у \setlist вище) — список літератури є окремим розділом зі своїм стандартним оформленням.
\usepackage[backend=biber,style=ieee]{biblatex}
% biblatex за замовчуванням додає розтяжний проміжок після «/ . : ? =» в адресах, і з вимкненими переносами рядок вирівнюється розрідженням самої адреси («https : / / www . scopus»). Рвати адресу й так дозволяє xurl, тож проміжок лише шкодить.
\setlength{\biburlbigskip}{0mu}
% Без додаткового інтервалу між описами: розділ має йти суцільним текстом, як решта звіту. Так само в шаблоні бакалаврської.
\setlength{\bibitemsep}{0pt}
\setlength{\bibparsep}{0pt}

% ---------------------------- Лістинги програмного коду ---------------------------- Використано minted (через Pygments), а не listings: listings не вміє коректно обробляти багатобайтові символи й «висмикує» кириличні коментарі з коду, зліплюючи їх в один рядок. Потребує прапорця -shell-escape.
\usepackage{xcolor}
\usepackage[newfloat]{minted}
% labelsep=period дає «Рис. 2.1. Назва» замість «Рис. 2: Назва» — формат НУЛП.
\usepackage[labelsep=period]{caption}

% baselinestretch=1 обов'язковий: без нього лістинг успадковує півторашку основного тексту, і довгий фрагмент коду перестає вміщатися у сторінку. Усередині [H]-флоата це не помилка збірки, а нескінченний цикл порожніх сторінок — lab1 так намолотив 161 тисячу сторінок.
\setminted{
  baselinestretch=1,
  fontsize=\footnotesize,
  linenos=true,
  frame=single,
  framesep=4pt,
  breaklines=true,
  breakanywhere=true,
  tabsize=4,
  numbersep=7pt,
  bgcolor=gray!8,
}
\setminted[python]{python3=true}

\SetupFloatingEnvironment{listing}{name=Лістинг,listname=Перелік лістингів}

% ---------------------------- Нумерація та оформлення заголовків ----------------------------
\usepackage{titlesec}

\setcounter{secnumdepth}{3}
\renewcommand{\thesection}{\arabic{section}}
\renewcommand{\thesubsection}{\arabic{section}.\arabic{subsection}}
\renewcommand{\thesubsubsection}{\arabic{section}.\arabic{subsection}.\arabic{subsubsection}}

% Розділ по центру, 16 пт, великими літерами; слово «РОЗДІЛ» опущене, бо звіт не дисертація. Номер лишається: \numberwithin прив'язує до нього нумерацію рисунків («Рис. 2.1»), і без нумерованих розділів вона б зламалася. \MakeUppercase четвертим аргументом отримує текст заголовка, тому в джерелі заголовки пишуться звичайним регістром.
\titleformat{\section}[block]
  {\centering\bfseries\fontsize{16}{19.2}\selectfont}
  {\thesection.}{0.3em}{\MakeUppercase}
% Підрозділи — зліва з абзацного відступу, звичайним регістром.
\titleformat{\subsection}[block]{\bfseries}{\thesubsection.}{1em}{}
\titleformat{\subsubsection}[block]{\bfseries\itshape}{\thesubsubsection.}{1em}{}

% Незіркові \titlespacing — обов'язкові. Без них titlesec застосовує зірковий варіант, який прибирає абзацний відступ у першого абзацу після заголовка й перекриває indentfirst. Значення — типові для report.
\titlespacing{\section}{0pt}{20pt}{6pt}
\titlespacing{\subsection}{1.27cm}{6pt}{0pt}
\titlespacing{\subsubsection}{1.27cm}{6pt}{0pt}

\setlength{\parindent}{1.27cm}
\setlength{\parskip}{0pt}
\renewcommand{\arraystretch}{1.3}

% Переноси слів вимкнені — так у шаблоні НУЛП. Натомість абзац розганяється міжслівними пробілами, тому tolerance і emergencystretch підняті: без них рядки виходили б за праве поле.
\hyphenpenalty=10000
\exhyphenpenalty=10000
\tolerance=9999
\setlength{\emergencystretch}{3em}

% ---------------------------- Підписи до рисунків і таблиць ----------------------------
\captionsetup[figure]{name=Рис.,justification=centering,font={bf,it}}
\captionsetup[table]{name=Таблиця,justification=raggedleft,singlelinecheck=false,labelfont=bf,textfont=normalfont}

% Рисунки, таблиці й формули нумеруються в межах розділу: «Рис. 2.1», а не «Рис. 5». Так у шаблоні НУЛП, і так само виглядали приклади на консультації.
\numberwithin{equation}{section}
\numberwithin{figure}{section}
\numberwithin{table}{section}

% ---------------------------- Колонтитули ----------------------------
\usepackage{fancyhdr}
\pagestyle{fancy}
\fancyhf{}
\fancyfoot[R]{\thepage}
\renewcommand{\headrulewidth}{0pt}
\fancypagestyle{plain}{%
  \fancyhf{}
  \fancyfoot[R]{\thepage}
  \renewcommand{\headrulewidth}{0pt}
}

\input{../../../../titlepage}
: \renewcommand{\WorkKind}{ПРАКТИЧНА РОБОТА}.
\newcommand{\WorkKind}{ЛАБОРАТОРНА РОБОТА}

\newcommand{\MakeLabTitle}[2]{%
\begin{titlepage}
% Титулка — не суцільний текст, а фіксована верстка з явних \\[...] проміжків,
% тож міжрядковий інтервал основного тексту (\setstretch{1.5} у preamble.tex)
% до неї не застосовний: він розтягував кожен із ~15 рядків на 20 %, і блок
% «Виконав» разом із «Львів --- рік» переповнював сторінку. Інтервал локальний
% для оточення titlepage, тому на решту документа це не впливає.
\setstretch{1}
\begin{center}
  {\large\bfseries МІНІСТЕРСТВО ОСВІТИ І НАУКИ УКРАЇНИ}\\[6pt]
  {\large\bfseries Національний університет «Львівська політехніка»}\\[6pt]
  {\large\bfseries Інститут комп'ютерних наук та інформаційних технологій}\\[6pt]
  {\large\bfseries Кафедра «Системи штучного інтелекту»}\\[1.2cm]

  \includegraphics[width=5cm]{../../../../logo}\\[1.2cm]

  {\Large\bfseries \WorkKind{} №#1}\\[0.5cm]
  з дисципліни\\[4pt]
  «\SubjectName»\\[0.6cm]
  на тему\\[4pt]
  \textit{#2}
\end{center}

% Відступ до блоку «Виконав» гнучкий, а не фіксований (було 2 cm). Фіксований
% разом із довгим родом роботи чи темою на два рядки переповнював сторінку, і
% «Львів --- рік» зі свого \vfill з'їжджав на окрему другу сторінку. \vfill
% стискається до нуля, коли місця бракує, тож переповнення стало неможливим.
\vfill

\begin{flushright}
  {\bfseries Виконав:}\\
  студент групи \StudentGroup\\
  \StudentName\\[1cm]
  {\bfseries Викладач:}\\
  \LecturerName\\
\end{flushright}

\vfill
\begin{center}
  Львів --- \AcademicYear
\end{center}
\end{titlepage}
\clearpage
}


% Subject metadata, read by the shared title page (titlepage.tex).

\newcommand{\SubjectName}{Методологія та процеси моделювання у штучному інтелекті}
% The lecturer's position and degree are not known yet; the other subjects put them on a line above the name, separated by `\\`.
\newcommand{\LecturerName}{Косаревич Р.Я.}
\newcommand{\StudentName}{Галич Максим}
\newcommand{\StudentGroup}{КНСШ-13}
\newcommand{\AcademicYear}{2026}

% Figures are generated artifacts, so they live in .cache/ rather than beside the text (see ARTIFACTS in src/mk1/config.py).
\graphicspath{{../../../../.cache/reports-artifacts/1th term 2026 Autumn/Методологія та процеси моделювання у штучному інтелекті/}}
% A centred column of fixed width: the wide tables below need their headers to wrap, which a plain c column cannot do.
\newcolumntype{C}[1]{>{\centering\arraybackslash}p{#1}}

\begin{document}

\MakeLabTitle{3}{Оптимізація гіперпараметрів, вкладена крос-валідація та дослідження витоку даних}

\textbf{Мета роботи.} Навчитися застосовувати методи оптимізації гіперпараметрів та крос-валідації; дослідити вплив вибору алгоритму пошуку на обчислювальну складність і якість моделей; опанувати методологію запобігання витоку даних та застосування вкладеної крос-валідації (Nested Cross-Validation).

Усі обчислення виконано мовою Python із бібліотеками scikit-learn [\ref{ref:sklearn}], imbalanced-learn та Optuna [\ref{ref:optuna}]. Код продовжує пакет \texttt{mk1} з лабораторних робіт №1--2. Пошук гіперпараметрів описано в модулі \texttt{tuning.py}, дослідження витоку даних --- у \texttt{leakage.py}, вкладену крос-валідацію та підсумкове порівняння --- у \texttt{nested.py}. Сценарій \texttt{lab3/run.py} лише запускає етапи по черзі й перевіряє, що кожен записав свої рисунки та метрики. Увесь експеримент відтворюється однією командою з фіксованим зерном генератора випадкових чисел.

\section{Вихідні дані та базова модель}
\label{sec:data}

Використано набір даних і розбиття з лабораторних робіт №1--2: \textbf{Diabetes 130-US Hospitals for Years 1999--2008} [\ref{ref:strack}], 69\,987 незалежних пацієнтів після очищення. Цільова змінна --- повторна госпіталізація протягом 30 днів після виписки (бінарна класифікація), клас 1 становить 8,98~\% спостережень. Навчальна вибірка містить 55\,989 пацієнтів, тестова --- 13\,998. Усі дослідження виконано крос-валідацією на навчальній вибірці. Тестову вибірку використано лише для контрольної перевірки висновків у розділах~\ref{sec:leakage} і~\ref{sec:summary}, і жодного вибору на ній не зроблено.

\subsection{Вибір моделі}

З моделей лабораторної роботи №2 обрано \textbf{гістограмний градієнтний бустинг} [\ref{ref:lightgbm}] на даних сценарію Б (повне очищення, Target Encoding високопотужних категоріальних ознак, доменні та агреговані ознаки). Причини вибору:

\begin{itemize}
\item у лабораторній роботі №2 він мав найвищу якість ранжування (PR-AUC 0,179--0,180 на крос-валідації);
\item він має кілька гіперпараметрів, що взаємодіють: крок навчання, розмір дерева, мінімальний розмір листка та регуляризація. Оптимальне значення одного залежить від інших, тому їх доводиться шукати разом;
\item він схильний до перенавчання: за великих дерев і великого кроку навчання ансамбль запам'ятовує навчальні записи;
\item він навчається швидко. Значення ознак групуються в 255 інтервалів, тому одне навчання на 37 тисячах записів триває кілька секунд, а вся робота потребує понад три тисячі навчань.
\end{itemize}

Випадковий ліс мав близьку якість, але його 300 дерев навчаються повільніше, а єдиний гіперпараметр, що суттєво впливав на якість (мінімальний розмір листка), уже досліджено в лабораторній роботі №2.

\subsection{Метрика якості}

Метрика якості --- \textbf{PR-AUC} (площа під кривою точність-повнота, у scikit-learn --- \texttt{average\_precision}), основна метрика лабораторних робіт №1--2. Модель видає ймовірність повторної госпіталізації, яку використовують для ранжування пацієнтів за ризиком, тож метрика має оцінювати саме ранжування, без вибору порога. За частки класу 1 близько 9~\% PR-AUC чутливіша до якості ранжування позитивного класу, ніж ROC-AUC [\ref{ref:saito}], яку обчислено як допоміжну метрику. Випадковий класифікатор має PR-AUC, що дорівнює частці класу 1 (0,090), і ROC-AUC 0,5. Accuracy непридатна: правило «завжди клас 0» дає 0,910.

\subsection{Базова модель}

Базова модель --- бустинг із гіперпараметрами лабораторної роботи №2: крок навчання 0,05, до 31 листка в дереві, щонайменше 20 записів у листку, L2-регуляризація 1, 250 дерев. Її оцінено двома схемами крос-валідації, з якими далі порівнюються результати оптимізації (табл.~\ref{tab:baseline}).

\begin{table}[H]
\caption{Якість базової моделі}
\label{tab:baseline}
\begin{tabularx}{\linewidth}{|X|c|c|c|}
\hline
\textbf{Схема оцінювання} & \textbf{PR-AUC} & \textbf{ROC-AUC} & \textbf{Час, с} \\
\hline
3-блокова крос-валідація (як у пошуку) & 0,1713 $\pm$ 0,0023 & 0,6491 & 13,7 \\
\hline
5-блокова крос-валідація (як зовнішня петля) & 0,1790 $\pm$ 0,0052 & 0,6550 & 21,3 \\
\hline
\end{tabularx}
\end{table}

5-блокова оцінка збігається з результатом лабораторної роботи №2. 3-блокова нижча на 0,008, хоча модель та сама: у 3-блоковій схемі модель навчається на двох третинах навчальної вибірки (37\,326 записів), а в 5-блоковій --- на чотирьох п'ятих (близько 44\,790 записів). Отже, бустинг ще не досяг насичення за обсягом даних, і абсолютні значення метрики в різних схемах не можна порівнювати напряму. Тому далі ефект оптимізації вимірюється як різниця з базовою моделлю в тій самій схемі.

\section{Простір гіперпараметрів і методи пошуку}
\label{sec:space}

\subsection{Простір пошуку}

Оптимізовано чотири гіперпараметри бустингу, що визначають складність моделі та силу регуляризації (табл.~\ref{tab:space}).

\begin{table}[H]
\caption{Простір гіперпараметрів}
\label{tab:space}
\small
\renewcommand{\arraystretch}{1.15}
\begin{tabularx}{\linewidth}{|p{3.4cm}|X|p{2.7cm}|p{3.0cm}|c|}
\hline
\textbf{Гіперпараметр} & \textbf{Зміст} & \textbf{Значення сітки} & \textbf{Розподіл для Random і Bayesian} & \textbf{Лаб. 2} \\
\hline
\texttt{learning\_rate} & крок навчання: частка, з якою прогноз нового дерева додається до ансамблю & 0,02; 0,063; 0,2 & логарифмічно рівномірний на $[0{,}02;\,0{,}2]$ & 0,05 \\
\hline
\texttt{max\_leaf\_nodes} & найбільша кількість листків у дереві, тобто складність одного дерева & 8; 32; 128 & цілий, логарифмічно рівномірний на $[8;\,128]$ & 31 \\
\hline
\texttt{min\_samples\_leaf} & найменша кількість записів у листку & 20; 100; 500 & цілий, логарифмічно рівномірний на $[20;\,500]$ & 20 \\
\hline
\texttt{l2\_regularization} & L2-штраф на значення листків & 0,1; 1; 10 & логарифмічно рівномірний на $[0{,}1;\,10]$ & 1 \\
\hline
\end{tabularx}
\end{table}

Для кожного гіперпараметра задано інтервал. Сітка містить його межі та середнє геометричне, тобто $3^4 = 81$ комбінацію. Випадковий і байєсівський пошук обирають точки з тих самих інтервалів, тож усі три методи досліджують одну область і відрізняються лише способом вибору точок у ній. Шкала всюди логарифмічна: вплив цих гіперпараметрів мультиплікативний, і зміна кроку навчання з 0,02 до 0,04 важить стільки ж, скільки з 0,1 до 0,2. Для цілих гіперпараметрів у scikit-learn немає логарифмічного розподілу (\texttt{randint} рівномірний на лінійній шкалі й обирав би переважно найбільші та найповільніші дерева), тому його реалізовано окремим класом \texttt{LogUniformInt}. Optuna для тих самих параметрів використовує \texttt{suggest\_int(log=True)}.

Кількість дерев зафіксовано на значенні 250, як у лабораторній роботі №2, а ранню зупинку вимкнено. Крок навчання і кількість дерев взаємозамінні: менший крок потребує більше дерев. Якщо оптимізувати обидва параметри, час навчання різних кандидатів відрізнявся б у десятки разів, і порівняння методів за часом втратило б сенс. Рання зупинка відділяє частину навчальних записів для внутрішньої перевірки, тож кандидати навчалися б на різних даних.

\subsection{Методи пошуку}

\begin{itemize}
\item \textbf{Повний перебір} (\texttt{GridSearchCV}) оцінює всі 81 комбінацію сітки.
\item \textbf{Випадковий пошук} (\texttt{RandomizedSearchCV}) [\ref{ref:bergstra}] оцінює $N = 20$ та $N = 50$ точок, вибраних незалежно з розподілів табл.~\ref{tab:space}. Обидва запуски використовують те саме зерно, а точки обираються послідовно одним генератором, тож 20 точок першого запуску збігаються з першими 20 точками другого. Різниця між ними показує внесок саме додаткових 30 точок.
\item \textbf{Байєсівська оптимізація} [\ref{ref:snoek}] реалізована бібліотекою Optuna з алгоритмом TPE (Tree-structured Parzen Estimator) [\ref{ref:tpe}], $N = 20$ спроб. TPE ділить уже оцінені точки на кращі та гірші й обирає наступну там, де відношення щільностей цих двох груп найбільше. Перші 5 спроб випадкові: моделі потрібні результати, на яких вона навчиться. Типове значення Optuna (10) витратило б на випадкові точки половину бюджету в 20 спроб. Використано багатовимірний TPE, що моделює гіперпараметри спільно: крок навчання і розмір дерева мають сенс лише разом.
\end{itemize}

Повний перебір і випадковий пошук виконують класи scikit-learn з однаковими параметрами оцінювання (лістинг~\ref{lst:sklearn}), байєсівську оптимізацію --- цільова функція Optuna, що оцінює кандидата тією самою крос-валідацією (лістинг~\ref{lst:bayes}). Кожен метод повертає історію всіх оцінених кандидатів, з якої обирається кандидат із найвищою середньою PR-AUC. Пошук не перенавчає остаточну модель (\texttt{refit=False}): час пошуку --- це час вибору гіперпараметрів.

\inputminted[firstline=244,lastline=261]{python}{../src/mk1/tuning.py}
\captionof{listing}{Повний перебір і випадковий пошук (\texttt{src/mk1/tuning.py})}
\label{lst:sklearn}

\inputminted[firstline=309,lastline=325]{python}{../src/mk1/tuning.py}
\captionof{listing}{Байєсівська оптимізація з Optuna (\texttt{src/mk1/tuning.py})}
\label{lst:bayes}

\subsection{Оцінювання кандидатів і вимірювання часу}

Кожного кандидата оцінено стратифікованою 3-блоковою крос-валідацією на всій навчальній вибірці. Блоки однакові для всіх методів, тож відмінності між методами зумовлені лише вибором кандидатів. На кожному блоці заново навчається весь конвеєр, разом із Target Encoding і масштабуванням, а не лише класифікатор. Ця ж 3-блокова схема є внутрішньою петлею вкладеної крос-валідації в розділі~\ref{sec:nested}, тому найкраща оцінка пошуку та внутрішня оцінка вкладеної схеми виміряні однаково.

Час виміряно двома способами:

\begin{itemize}
\item \textbf{тривалість пошуку} --- час від початку до кінця пошуку на комп'ютері з чотирма ядрами;
\item \textbf{сумарний обчислювальний час} --- сума тривалостей усіх навчань і оцінювань кандидатів на всіх блоках. Він не залежить від того, скільки навчань виконувалося паралельно, і пропорційний обсягу роботи методу.
\end{itemize}

Ці величини відрізняються через паралелізм. \texttt{GridSearchCV} і \texttt{RandomizedSearchCV} знають усіх кандидатів наперед і навчають їх у чотирьох процесах одночасно. Байєсівська оптимізація обирає наступну точку за результатами всіх попередніх, тому спроби виконуються одна за одною, і паралельно навчаються лише три блоки однієї спроби. Отже, байєсівський пошук використовує три ядра з чотирьох. Це властивість методу, а не реалізації, і її враховано у висновках.

\section{Дослідження 1: порівняння методів оптимізації}
\label{sec:tuning}

Кожен метод виконано на всій навчальній вибірці. Результати наведено в табл.~\ref{tab:tuning}, обрані гіперпараметри --- в табл.~\ref{tab:best}. Приріст обчислено відносно базової моделі в тій самій 3-блоковій схемі (PR-AUC 0,1713).

\begin{table}[H]
\caption{Порівняння методів пошуку на 3-блоковій крос-валідації}
\label{tab:tuning}
\small
\renewcommand{\arraystretch}{1.15}
\begin{tabularx}{\linewidth}{|X|C{1.5cm}|C{1.9cm}|C{1.9cm}|C{2.8cm}|C{1.6cm}|C{1.5cm}|}
\hline
\textbf{Метод} & \textbf{Канди\-датів} & \textbf{Трива\-лість, с} & \textbf{Обчисл. час, с} & \textbf{Найкраща PR-AUC} & \textbf{При\-ріст} & \textbf{ROC-AUC} \\
\hline
Без оптимізації (лаб. 2) & 1 & 14 & 30 & 0,1713 $\pm$ 0,0023 & --- & 0,6491 \\
\hline
Grid Search & 81 & 766 & 2972 & 0,1783 $\pm$ 0,0048 & +0,0070 & 0,6562 \\
\hline
Random Search, $N = 20$ & 20 & 182 & 685 & 0,1770 $\pm$ 0,0059 & +0,0057 & 0,6564 \\
\hline
Random Search, $N = 50$ & 50 & 459 & 1761 & 0,1770 $\pm$ 0,0059 & +0,0057 & 0,6564 \\
\hline
Bayesian (TPE), $N = 20$ & 20 & 220 & 612 & 0,1779 $\pm$ 0,0045 & +0,0066 & 0,6565 \\
\hline
\end{tabularx}
\end{table}

\begin{table}[H]
\caption{Гіперпараметри, обрані кожним методом}
\label{tab:best}
\small
\renewcommand{\arraystretch}{1.15}
\begin{tabularx}{\linewidth}{|X|C{2.4cm}|C{2.4cm}|C{2.6cm}|C{2.4cm}|}
\hline
\textbf{Метод} & \textbf{Крок навчання} & \textbf{Листків у дереві, до} & \textbf{Записів у листку, від} & \textbf{L2-регуля\-ризація} \\
\hline
Без оптимізації (лаб. 2) & 0,05 & 31 & 20 & 1 \\
\hline
Grid Search & 0,02 & 32 & 100 & 10 \\
\hline
Random Search, $N = 20$ і $N = 50$ & 0,033 & 13 & 36 & 4,62 \\
\hline
Bayesian (TPE), $N = 20$ & 0,020 & 18 & 110 & 2,39 \\
\hline
\end{tabularx}
\end{table}

\begin{figure}[H]
  \centering
  \includegraphics[width=\linewidth]{lab3/tuning_progress}
  \caption{Найкраща PR-AUC, знайдена до певного моменту, залежно від сумарного обчислювального часу. Точки --- окремі кандидати в порядку їх оцінювання; кандидати з PR-AUC нижче 0,148 за межами графіка}
  \label{fig:progress}
\end{figure}

\begin{figure}[H]
  \centering
  \includegraphics[width=\linewidth]{lab3/tuning_space}
  \caption{Точки, які оцінив кожен метод, у площині «крок навчання --- кількість листків», колір --- PR-AUC. У повному переборі кожна точка площини поєднує 9 кандидатів з різними значеннями двох інших гіперпараметрів, зверху показано найкращого з них}
  \label{fig:space}
\end{figure}

\textbf{Усі методи покращили модель приблизно однаково.} Приріст становить 0,006--0,007 PR-AUC, тобто 3--4~\% від базового значення. Найкращі оцінки методів відрізняються не більше ніж на 0,0013, а стандартне відхилення найкращого кандидата між блоками становить 0,005. Отже, за якістю методи в межах похибки еквівалентні, і різниця між ними --- це насамперед вартість. Поверхня якості пласка поблизу оптимуму: у кожного методу від 6 до 8 кандидатів відстають від найкращого менше ніж на 0,002 PR-AUC.

Усі методи обрали менший крок навчання, ніж у лабораторній роботі №2 (0,02--0,033 замість 0,05), і сильнішу регуляризацію: більший мінімальний розмір листка (36--110 записів замість 20) та більший L2-штраф (2,4--10 замість 1). Отже, базова модель зі 250 деревами та кроком 0,05 дещо перенавчена. Grid Search і TPE обрали крок 0,02, тобто нижню межу інтервалу пошуку. Оптимум може лежати нижче, але менший крок потребує більше дерев і, відповідно, більше часу на кожного кандидата. Це обмеження заданого простору пошуку, а не методів.

\textbf{Повний перебір} знайшов найвищу оцінку (0,1783), але тривав у 3,5 раза довше за байєсівську оптимізацію й виконав у 4,9 раза більше обчислень заради різниці в 0,0004, що значно менша за похибку. Сітка має 81 точку, але лише по три різні значення кожного гіперпараметра (рис.~\ref{fig:space}, ліворуч): якщо важливий лише один параметр, 81 оцінка дає про нього стільки ж інформації, скільки три [\ref{ref:bergstra}]. На рис.~\ref{fig:progress} крива повного перебору рано досягає високого рівня лише тому, що перші кандидати сітки мають найменший крок навчання, який і виявився найкращим. За іншого порядку значень перебір довго оцінював би гірші точки. Вартість повного перебору зростає експоненційно: п'ятий гіперпараметр із трьома значеннями збільшив би її втричі, до 243 кандидатів.

\textbf{Випадковий пошук} знайшов свого найкращого кандидата вже четвертою точкою, приблизно за 2 хвилини обчислювального часу. Подальші 16 точок ($N = 20$), а потім ще 30 ($N = 50$) його не покращили: лише одна з нових точок дала майже такий самий результат. Збільшення бюджету в 2,5 раза не дало нічого, бо пласку область поблизу оптимуму випадкові точки вже знайшли, а щоб потрапити ще ближче до межі інтервалу, де лежить оптимум, потрібне везіння, а не кількість спроб. Випадкові точки рівномірно вкривають увесь простір (рис.~\ref{fig:space}, у центрі), і 50 різних значень кожного гіперпараметра дали б перевагу над сіткою, якби якість залежала від одного з них різко.

\textbf{Байєсівська оптимізація} покращувала результат протягом усього бюджету: 0,1770 після другої спроби, 0,1774 після шостої, 0,1778 після шістнадцятої і 0,1779 на останній, двадцятій. TPE швидко зосередився на області малого кроку навчання й невеликих дерев (рис.~\ref{fig:space}, праворуч): 14 з 20 його кандидатів кращі за базову модель, проти 16 з 50 у випадкового пошуку та 32 з 81 у сітці. Медіана PR-AUC його кандидатів --- 0,1742, у випадкового пошуку --- 0,1684, у сітки --- 0,1671. TPE також витратив найменше обчислювального часу (612~с), бо обирав невеликі дерева, які навчаються швидше. Проте за тривалістю він повільніший за випадковий пошук з тим самим бюджетом (220~с проти 182~с): спроби виконуються послідовно, і з чотирьох ядер використовуються лише три. Відношення обчислювального часу до тривалості становить 2,8 для TPE проти 3,8--3,9 для інших методів.

\section{Дослідження 2: вплив витоку даних}
\label{sec:leakage}

\subsection{Варіанти конвеєра}

Витік даних (data leakage) --- це ситуація, коли модель під час навчання чи вибору отримує інформацію, якої не матиме під час застосування до нових записів [\ref{ref:kaufman}]. Під час крос-валідації найчастіше джерело витоку --- кроки підготовки даних, навчені на всій вибірці ще до її поділу на блоки. Тоді валідаційний блок уже «бачили» кроки, що його перетворюють, і оцінка якості більше не відповідає ситуації з новим пацієнтом.

Порівняно чотири варіанти (лістинг~\ref{lst:leakage}). У всіх варіантах ті самі записи, та сама модель (бустинг із гіперпараметрами лабораторної роботи №2) і ті самі 5 блоків крос-валідації. Відрізняється лише те, які кроки навчаються на всій навчальній вибірці до поділу на блоки.

\begin{enumerate}
\item \textbf{Коректний конвеєр.} Увесь конвеєр, від інженерії ознак до класифікатора, передано в крос-валідацію. На кожному навчальному блоці всі кроки навчаються заново, а до валідаційного блоку лише застосовуються.
\item \textbf{Препроцесинг до розбиття (\texttt{fit\_transform}).} Усі кроки підготовки навчено на всій навчальній вибірці методом \texttt{fit\_transform}, а в крос-валідацію передано лише класифікатор. Групові середні агрегатних ознак, моди для пропусків і параметри масштабування обчислено разом із валідаційними записами. Target Encoding у \texttt{fit\_transform} бібліотеки scikit-learn навчається з перехресним навчанням: кожен рядок кодується середніми, обчисленими на інших блоках, тобто без власної мітки.
\item \textbf{Препроцесинг до розбиття (\texttt{fit}, потім \texttt{transform}).} Ті самі кроки навчено методом \texttt{fit} на всій навчальній вибірці, яку потім перетворено методом \texttt{transform}. Це звичний рецепт для масштабування, але для Target Encoding він означає, що кожен рядок закодовано середнім значенням цілі, обчисленим разом із його власною міткою. Саме цей варіант є некоректним конвеєром, якого вимагає завдання.
\item \textbf{Препроцесинг і SMOTE до розбиття.} Варіант 2, після якого всю навчальну вибірку доповнено синтетичними записами класу 1 методом SMOTE [\ref{ref:smote}] до співвідношення 1~:~1, і лише потім поділено на блоки.
\end{enumerate}

\inputminted[firstline=64,lastline=92]{python}{../src/mk1/leakage.py}
\captionof{listing}{Підготовка даних у варіантах із витоком (\texttt{src/mk1/leakage.py})}
\label{lst:leakage}

Кожен варіант також навчено на всій навчальній вибірці й оцінено на тестовій. Тестова вибірка не брала участі в навчанні жодного кроку жодного варіанта, тож її оцінка чесна. Різниця між оцінкою крос-валідації та тестовою показує, наскільки витік завищив оцінку.

\subsection{Результати}

Результати наведено в табл.~\ref{tab:leakage} та на рис.~\ref{fig:leakage}. Стовпець «Завищення» --- різниця між оцінкою крос-валідації варіанта та коректного конвеєра. Стовпець «Зміна на тесті» --- така сама різниця на тестовій вибірці, тобто реальна зміна якості моделі.

\begin{table}[H]
\caption{Оцінки якості за різної підготовки даних відносно поділу на блоки}
\label{tab:leakage}
\small
\renewcommand{\arraystretch}{1.15}
\begin{tabularx}{\linewidth}{|X|c|c|c|c|c|c|}
\hline
& \multicolumn{4}{c|}{\textbf{PR-AUC}} & \multicolumn{2}{c|}{\textbf{ROC-AUC}} \\
\cline{2-7}
\textbf{Варіант} & \textbf{CV} & \textbf{Завищення} & \textbf{Тест} & \textbf{Зміна на тесті} & \textbf{CV} & \textbf{Тест} \\
\hline
1. Коректний конвеєр & 0,179 $\pm$ 0,005 & --- & 0,189 & --- & 0,655 & 0,654 \\
\hline
2. Препроцесинг до розбиття (\texttt{fit\_transform}) & 0,182 $\pm$ 0,005 & +0,003 & 0,189 & 0,000 & 0,655 & 0,654 \\
\hline
3. Препроцесинг до розбиття (\texttt{fit} + \texttt{transform}) & 0,230 $\pm$ 0,008 & +0,051 & 0,161 & $-$0,028 & 0,711 & 0,633 \\
\hline
4. Препроцесинг і SMOTE до розбиття & 0,885 $\pm$ 0,052 & +0,706 & 0,101 & $-$0,088 & 0,829 & 0,516 \\
\hline
\end{tabularx}
\end{table}

\begin{figure}[H]
  \centering
  \includegraphics[width=\linewidth]{lab3/leakage}
  \caption{Оцінка крос-валідації (середнє $\pm$ стандартне відхилення за 5 блоками) і тестова оцінка кожного варіанта. Пунктир --- випадковий класифікатор на реальних даних. У варіанті 4 валідаційні блоки наполовину складаються з синтетичних записів, і для них PR-AUC випадкового класифікатора становить 0,5}
  \label{fig:leakage}
\end{figure}

\textbf{Некоректний конвеєр (варіант 3) завищив PR-AUC на 0,051, тобто на 28~\%}: крос-валідація показує 0,230 замість 0,179, а ROC-AUC --- 0,711 замість 0,655. Розкид між блоками (0,008) у шість разів менший за завищення, тож за самими результатами крос-валідації помилку не видно: дослідник просто вирішив би, що модель стала кращою. Насправді модель стала гіршою. На тестовій вибірці її PR-AUC 0,161 проти 0,189 у коректного конвеєра, а ROC-AUC 0,633 проти 0,654.

Механізм витоку пов'язаний із високопотужними ознаками. Коди діагнозів \texttt{diag\_1}--\texttt{diag\_3} мають 671--731 різних значень у навчальній вибірці, і більше половини з них трапляються не більше ніж у 10 пацієнтів (91--141 код --- лише в одного). Target Encoding замінює категорію середнім значенням цілі за її записами. Коли це середнє обчислено разом із власною міткою рядка, для рідкісного коду воно майже дорівнює цій мітці: пацієнт з унікальним кодом і повторною госпіталізацією отримує кодування, близьке до 1. Згладжування в \texttt{TargetEncoder} лише наближає таке кодування до загального середнього, але не прибирає залежність від власної мітки. Бустинг знаходить цю залежність, бо вона найсильніша в даних, і на валідаційних блоках вона так само працює: їхні записи закодовано з урахуванням їхніх власних міток. Новий пацієнт такої підказки не має, тож модель, що покладається на неї, на тесті прогнозує гірше за модель, навчену на чесному кодуванні.

\textbf{Варіант 2 майже не завищує оцінку} (+0,003, менше за розкид між блоками), а остаточна модель у ньому та сама, що й у коректному конвеєрі, тому тестові оцінки збігаються. Причина в тому, що \texttt{TargetEncoder.fit\_transform} у scikit-learn кодує кожен рядок середніми, обчисленими на інших блоках, тобто саме для того, щоб власна мітка не потрапила в кодування. Решта кроків, навчених до розбиття (групові середні, моди, параметри масштабування), цілі не використовують, і для деревоподібної моделі їхній витік несуттєвий. Отже, небезпечний не будь-який препроцесинг до розбиття, а той, що навчається на цілі. Проте покладатися на цей захист не варто: він працює лише за виклику \texttt{fit\_transform}, і варіант 3 відрізняється від варіанта 2 тільки тим, що \texttt{fit} і \texttt{transform} викликано окремо.

\textbf{SMOTE до розбиття (варіант 4) дає найбільше завищення}: PR-AUC 0,885 на крос-валідації проти 0,101 на тесті, а ROC-AUC 0,829 проти 0,516, що майже не відрізняється від випадкового ранжування. Тут поєднано два ефекти. По-перше, валідаційні блоки наполовину складаються з синтетичних записів класу 1, тож частка класу 1 у них 50~\% замість 9~\%, і PR-AUC навіть випадкового класифікатора зростає до 0,5. По-друге, кожен синтетичний запис лежить на відрізку між двома реальними записами класу 1, які здебільшого опинилися в навчальних блоках. Модель розпізнає синтетичні записи як «майже копії» навчальних. У лабораторній роботі №2 показано, що деревоподібні моделі розпізнають синтетичні записи майже безпомилково за дробовими значеннями закодованих ознак. Остаточна модель варіанта 4, навчена на всій навчальній вибірці, збігається з моделлю сценарію В (SMOTE) лабораторної роботи №2, і її ROC-AUC на тесті та сама: 0,516. Там передискретизацію виконано коректно, лише на навчальних блоках, і крос-валідація чесно показала падіння якості. Отже, SMOTE шкодить цій моделі і без витоку, а витік не просто завищує оцінку, а приховує цю шкоду, показуючи майже ідеальну якість.

Висновок дослідження: оцінка крос-валідації коректна лише тоді, коли всі кроки, що навчаються на даних, особливо ті, що використовують ціль (Target Encoding, передискретизація, відбір ознак), знаходяться всередині конвеєра й навчаються заново на кожному блоці. Розкид між блоками не виявляє витоку. Виявити його можна лише порівнянням з оцінкою на даних, які жоден крок не бачив.

\section{Дослідження 3: вкладена крос-валідація}
\label{sec:nested}

\subsection{Схема}

Найкраща оцінка пошуку з розділу~\ref{sec:tuning} зміщена вгору. Пошук оцінює десятки кандидатів на тих самих блоках і обирає максимум. Кожна оцінка містить випадкову похибку, тож максимум відбирає не лише кращий набір гіперпараметрів, а й найвдалішу похибку [\ref{ref:varma}, \ref{ref:cawley}]. Чим більше кандидатів, тим сильніший цей ефект. Щоб отримати незміщену оцінку, вибір гіперпараметрів і оцінювання мають використовувати різні дані. Це забезпечує \textbf{вкладена крос-валідація} (лістинг~\ref{lst:nested}):

\begin{itemize}
\item \textbf{зовнішня петля} ділить навчальну вибірку на $K = 5$ стратифікованих блоків. Кожен блок по черзі стає тестовим для оцінки узагальнювальної здатності, а решта чотири --- навчальними;
\item \textbf{внутрішня петля} на навчальних блоках зовнішньої петлі виконує весь пошук гіперпараметрів з 3-блоковою крос-валідацією, так само як у розділі~\ref{sec:tuning};
\item обрані гіперпараметри навчаються на всіх чотирьох навчальних блоках зовнішньої петлі й один раз оцінюються на п'ятому, якого пошук не бачив.
\end{itemize}

Середнє п'яти зовнішніх оцінок --- незміщена оцінка якості \emph{методу} пошуку, а не одного обраного кандидата: на кожному зовнішньому блоці метод може обрати інші гіперпараметри. Для кожного методу виконано окрему вкладену крос-валідацію з тими самими блоками.

\inputminted[firstline=91,lastline=107]{python}{../src/mk1/nested.py}
\captionof{listing}{Вкладена крос-валідація одного методу пошуку (\texttt{src/mk1/nested.py})}
\label{lst:nested}

Внутрішні оцінки отримано на моделях, навчених на двох третинах зовнішньої навчальної частини (близько 29,9 тис. записів), а зовнішні --- на всій частині (близько 44,8 тис.). Табл.~\ref{tab:baseline} показує, що вже ця різниця в обсязі даних змінює PR-AUC. Тому базову модель без оптимізації оцінено на тих самих внутрішніх і зовнішніх блоках. Її різниця між внутрішньою та зовнішньою оцінкою --- це ефект обсягу даних, а різниця методу понад неї --- оптимізм вибору.

\section{Підсумкове порівняння}
\label{sec:summary}

%% SUMMARY

\section{Висновки}

%% CONCLUSIONS

\section*{Список використаних джерел}

\begin{enumerate}[label={[\arabic*]},ref={\arabic*},leftmargin=*,labelindent=0pt]
\item \label{ref:strack} Strack B., DeShazo J. P., Gennings C. et al. Impact of HbA1c Measurement on Hospital Readmission Rates: Analysis of 70,000 Clinical Database Patient Records // BioMed Research International. 2014. Vol. 2014. Article ID 781670.
\item \label{ref:sklearn} Pedregosa F., Varoquaux G., Gramfort A. et al. Scikit-learn: Machine Learning in Python // Journal of Machine Learning Research. 2011. Vol. 12. P. 2825--2830.
\item \label{ref:optuna} Akiba T., Sano S., Yanase T., Ohta T., Koyama M. Optuna: A Next-generation Hyperparameter Optimization Framework // Proc. 25th ACM SIGKDD International Conference on Knowledge Discovery \& Data Mining. 2019. P. 2623--2631.
\item \label{ref:lightgbm} Ke G., Meng Q., Finley T. et al. LightGBM: A Highly Efficient Gradient Boosting Decision Tree // Advances in Neural Information Processing Systems 30. 2017. P. 3146--3154.
\item \label{ref:saito} Saito T., Rehmsmeier M. The Precision-Recall Plot Is More Informative than the ROC Plot When Evaluating Binary Classifiers on Imbalanced Datasets // PLoS ONE. 2015. Vol. 10, No. 3. e0118432.
\item \label{ref:bergstra} Bergstra J., Bengio Y. Random Search for Hyper-Parameter Optimization // Journal of Machine Learning Research. 2012. Vol. 13. P. 281--305.
\item \label{ref:tpe} Bergstra J., Bardenet R., Bengio Y., Kégl B. Algorithms for Hyper-Parameter Optimization // Advances in Neural Information Processing Systems 24. 2011. P. 2546--2554.
\item \label{ref:snoek} Snoek J., Larochelle H., Adams R. P. Practical Bayesian Optimization of Machine Learning Algorithms // Advances in Neural Information Processing Systems 25. 2012. P. 2951--2959.
\item \label{ref:kaufman} Kaufman S., Rosset S., Perlich C., Stitelman O. Leakage in Data Mining: Formulation, Detection, and Avoidance // ACM Transactions on Knowledge Discovery from Data. 2012. Vol. 6, No. 4. Article 15.
\item \label{ref:target} Micci-Barreca D. A Preprocessing Scheme for High-Cardinality Categorical Attributes in Classification and Prediction Problems // ACM SIGKDD Explorations Newsletter. 2001. Vol. 3, No. 1. P. 27--32.
\item \label{ref:smote} Chawla N. V., Bowyer K. W., Hall L. O., Kegelmeyer W. P. SMOTE: Synthetic Minority Over-sampling Technique // Journal of Artificial Intelligence Research. 2002. Vol. 16. P. 321--357.
\item \label{ref:varma} Varma S., Simon R. Bias in Error Estimation When Using Cross-Validation for Model Selection // BMC Bioinformatics. 2006. Vol. 7. Article 91.
\item \label{ref:cawley} Cawley G. C., Talbot N. L. C. On Over-fitting in Model Selection and Subsequent Selection Bias in Performance Evaluation // Journal of Machine Learning Research. 2010. Vol. 11. P. 2079--2107.
\end{enumerate}

\end{document}
